% Chapter 1 --- Cones (paper §2, milestone M1).
%
% Mirrors theories/cones/:
%   precone.v          paper Def 2.1, 2.2: precones
%   cone.v             paper Def 2.4: cones, norm, ω-completeness
%   basic_lemmas.v     supporting lemmas on cones
%   cone_cat.v         paper §2.4: the category Cones
%   examples_cone.v    canonical examples (⊤, ⊥)

%%%%%%%%%%%%%%%%%%%%%%%%%%%%%%%%%%%%%%%%%%%%%%%%%%%%%%%%%%%%%%%%%%%%%%%%
\chapter{Cones}\label{ch:cones}

This chapter formalises paper \S2 of Ehrhard and
Geoffroy~\cite{ehrhard-geoffroy-icones}: precones, cones, the
norm, $\omega$-completeness of the unit ball, and the category
$\Cones$. All material in this chapter lives under
\texttt{theories/cones/}.

Positive cones in the sense of Selinger~\cite{ehrhard-geoffroy-icones}
are algebraic structures with both a numerical and a domain-theoretic
flavour: an $\Rnn$-semimodule satisfying two cancellation/positivity
axioms (this is the \emph{precone} layer), augmented with a norm
$\cnorm{\cdot} : P \to \Rnn$ satisfying scalar-homogeneity,
zero-detection, sub-additivity, order-monotonicity, and the
$\omega$-completeness of the unit ball. They provide the basic
objects of the model of probabilistic computation developed in
\cite{ehrhard-geoffroy-icones}, where the algebraic and numerical
parts encode probabilistic combinations of states, and the
domain-theoretic part is what enables interpretation of recursion.
Compared with Tix~\cite{ehrhard-geoffroy-icones}, the order-theoretic
completeness is restricted to countable, norm-bounded
$\omega$-chains, which suffices for fixed-point reasoning and is
preserved by the monotone-convergence theorem of measure theory ---
a fact that becomes essential in M2 and M3.

%%%%%%%%%%%%%%%%%%%%%%%%%%%%%%%%%%%%%%%%%%%%%%%%%%%%%%%%%%%%%%%%%%%%%%%%
\section{Basic definitions}\label{sec:precones}

We begin with the purely algebraic layer.

\begin{definition}[Precone]\label{def:precone}
\rocqok\rocq{Icones.cones.precone.isPrecone}
An $\Rnn$-\emph{precone} is a set $P$ equipped with a constant
$0 \in P$, an addition $+ : P \times P \to P$ and a scalar
multiplication ${\cdot}: \Rnn \times P \to P$, making $P$ into a
commutative $\Rnn$-semimodule (associativity, commutativity, unit
$0$, distributivity, scalar associativity, $1 \cdot x = x$,
$r \cdot 0 = 0$, $0 \cdot x = 0$), and additionally satisfying:
\begin{description}
  \item[\textnormal{(Cancel)}] $\forall x, x_1, x_2 \in P,\;
    x_1 + x = x_2 + x \;\Rightarrow\; x_1 = x_2$;
  \item[\textnormal{(Pos)}] $\forall x_1, x_2 \in P,\;
    x_1 + x_2 = 0 \;\Rightarrow\; x_1 = 0 \text{ and } x_2 = 0$.
\end{description}
\end{definition}

The cancellation and positivity axioms together imply that the
relation $x_1 \cle x_2 \;\Leftrightarrow\; \exists z\in P,\,
x_2 = x_1 + z$ is a partial order on $P$, the \emph{cone order},
and that when $x_1 \cle x_2$ the difference $x_2 - x_1$ is uniquely
defined. The Rocq counterpart is $\rocq{Icones.cones.precone.precone_le}$,
with the antisymmetry/transitivity lemmas just below it in
\texttt{precone.v}. The norm-free precone layer is sufficient to talk
about the cone order and its arithmetic.

\begin{definition}[Cone]\label{def:cone}
\uses{def:precone}
\rocqok\rocq{Icones.cones.cone.isCone}
A \emph{cone} is a precone $P$ equipped with a function
$\cnorm{\cdot}_P : P \to \Rnn$, called the norm of $P$, satisfying
\begin{description}
  \item[\textnormal{(Normh)}]
    $\cnorm{\lambda x} = \lambda \cnorm{x}$ for all
    $\lambda \in \Rnn, x \in P$;
  \item[\textnormal{(Normz)}]
    $\cnorm{x} = 0 \;\Rightarrow\; x = 0$;
  \item[\textnormal{(Normt)}]
    $\cnorm{x_1 + x_2} \leq \cnorm{x_1} + \cnorm{x_2}$;
  \item[\textnormal{(Normp)}]
    $x_1 \cle x_2 \;\Rightarrow\; \cnorm{x_1} \leq \cnorm{x_2}$;
  \item[\textnormal{(Normc)}]
    Every $\omega$-chain $(x_n)_{n \in \N}$ in $P$ with
    $\cnorm{x_n} \leq 1$ for all $n$ has a least upper bound
    $x = \bigsup_{n} x_n \in P$ which moreover satisfies
    $\cnorm{x} \leq 1$.
\end{description}
We write $B_P = \{x \in P \mid \cnorm{x} \leq 1\}$ for the
\emph{unit ball} of $P$.
\end{definition}

The Rocq mixin packages (Normc) as a concrete operator
$\rocq{Icones.cones.cone.cone_sup_ball}$
together with three lemmas
($\rocq{Icones.cones.cone.cone_sup_ball_ub}$,
$\rocq{Icones.cones.cone.cone_sup_ball_lub}$,
$\rocq{Icones.cones.cone.cone_sup_ball_norm}$)
witnessing that it is an upper bound, the least one, and lives in the
unit ball. From (Normh) and (Normp) one immediately derives that
$\cnorm{0} = 0$ ($\rocq{Icones.cones.cone.cone_norm0}$), that
$\cnorm{x} \geq 0$ ($\rocq{Icones.cones.cone.cone_norm_ge0}$), and
that (Normz) can be repackaged as the biconditional
$\rocq{Icones.cones.cone.cone_normz_iff}$.

We use $\rocq{Icones.cones.cone.cnorm}$ as a short abbreviation for
$\rocq{Icones.cones.cone.cone_norm}$ throughout. Functions between
cones (or, more generally, between $\omega$-closed subsets of cones)
will be called \emph{increasing}, $\omega$-\emph{continuous},
\emph{linear}, or \emph{bounded} in the senses of paper Def.~2.2;
the corresponding Rocq predicates are
$\rocq{Icones.cones.basic_lemmas.is_increasing}$,
$\rocq{Icones.cones.basic_lemmas.is_omega_continuous}$, and
$\rocq{Icones.cones.basic_lemmas.is_linear}$.

%%%%%%%%%%%%%%%%%%%%%%%%%%%%%%%%%%%%%%%%%%%%%%%%%%%%%%%%%%%%%%%%%%%%%%%%
\section{The cone of finite measures (motivational example)}\label{sec:cone-fmeas}

Paper \S2.2 introduces the central example that drives the whole
development: for a measurable space $X$, the set $\FMeas(X)$ of finite
non-negative measures on $X$ is a cone. The algebraic operations are
defined pointwise on measurable sets ($(\mu_1 + \mu_2)(U) =
\mu_1(U) + \mu_2(U)$, $(\lambda \mu)(U) = \lambda \mu(U)$); the norm
is the total mass $\cnorm{\mu} = \mu(X)$; the cone order coincides
with the pointwise order; least upper bounds of norm-bounded
increasing sequences are computed pointwise on measurable sets and
exist by the monotone-convergence theorem. The proper treatment of
$\FMeas(X)$ as a cone (and later as a measurable cone) is the subject
of Chapter~\ref{ch:mcones}. We mention it here only as the prototype
which all the abstract conditions of Definition~\ref{def:cone} are
designed to capture.

%%%%%%%%%%%%%%%%%%%%%%%%%%%%%%%%%%%%%%%%%%%%%%%%%%%%%%%%%%%%%%%%%%%%%%%%
\section{Basic properties}\label{sec:cone-basic-lemmas}

The next four lemmas are the algebraic backbone of the theory. They
appear as paper Lemmas~2.8--2.11. The Rocq counterparts live in
$\rocq{Icones.cones.basic_lemmas}$.

\begin{lemma}[Inverse of a linear $\omega$-continuous bijection]
\label{lem:invf-linear-cont}
\uses{def:cone}
\rocqok\rocq{Icones.cones.basic_lemmas.invf_linear}
\rocq{Icones.cones.basic_lemmas.invf_omega_continuous}
Let $P$ and $Q$ be cones and let $f : P \to Q$ be a linear
$\omega$-continuous bijection. Then $f^{-1} : Q \to P$ is also
linear and $\omega$-continuous.
\end{lemma}

\begin{proof}
Linearity of $f^{-1}$ follows from injectivity of $f$: for
$y_1, y_2 \in Q$ both $f^{-1}(y_1 + y_2)$ and $f^{-1}(y_1) +
f^{-1}(y_2)$ have image $y_1 + y_2$ under $f$, so they agree;
similarly for scalar multiplication. Since $f^{-1}$ is linear it
is monotone. For $\omega$-continuity, let $(y_n)$ be an
$\omega$-chain in $B_Q$ with sup $y$. The chain $(f^{-1}(y_n))$
is increasing, bounded above by $f^{-1}(y)$, hence sits in a
ball of $P$ and has a sup $x \cle f^{-1}(y)$. Applying $f$ and
using its continuity gives $f(x) = \sup_n y_n = y$, so
$x = f^{-1}(y)$ by injectivity.
\end{proof}

\begin{lemma}[Addition and scaling are increasing and $\omega$-continuous]
\label{lem:add-scale-cont}
\uses{def:cone}
\rocqok\rocq{Icones.cones.basic_lemmas.addr_increasing}
\rocq{Icones.cones.basic_lemmas.addl_increasing}
\rocq{Icones.cones.basic_lemmas.scaler_increasing}
\rocq{Icones.cones.basic_lemmas.addr_omega_continuous}
\rocq{Icones.cones.basic_lemmas.addl_omega_continuous}
\rocq{Icones.cones.basic_lemmas.scaler_omega_continuous}
Let $P$ be a cone. The maps
$x \mapsto x + y$, $y \mapsto x + y$ and $x \mapsto \lambda \cdot x$
(for fixed $y \in P$, $x \in P$, $\lambda \in \Rnn$) are increasing
and $\omega$-continuous on $P$.
\end{lemma}

\begin{proof}
Monotonicity is by inspection of the cone order: if $x_1 \cle x_2$
with witness $x_2 = x_1 + z$, then $x_1 + y \cle x_2 + y$ via the
same $z$, and similarly for scaling. For $\omega$-continuity, one
shows directly that the chain sup commutes with $x \mapsto x + y$
by appealing to least-upper-bound uniqueness in the unit ball
(after a renormalisation: $(x_n + y)$ may exit $B_P$, so the proof
proceeds by rescaling to a chain in $B_P$, taking the sup there,
and rescaling back).
\end{proof}

\begin{lemma}[$\omega$-continuity of differences]
\label{lem:diff-omega-continuous}
\uses{lem:add-scale-cont}
\rocqok\rocq{Icones.cones.basic_lemmas.diff_omega_continuous}
Let $P$ and $Q$ be cones, $A \subseteq P$ an $\omega$-closed
subset, and $f, g : A \to Q$ functions with $f$ increasing, $g$
$\omega$-continuous, and $f(x) \cle g(x)$ for all $x \in A$. If
the pointwise difference $g - f : x \mapsto g(x) - f(x)$ is
increasing, then $g - f$ is $\omega$-continuous.
\end{lemma}

\begin{proof}
Let $(x_n)$ be a bounded increasing chain in $A$ with sup $x$.
For each $n$, $f(x_n) \cle f(x) \cle f(x) + (g(x_n) - f(x_n))$,
i.e.\ $g(x_n) \cle f(x) + (g(x_n) - f(x_n))$. The chain $n
\mapsto f(x) + (g(x_n) - f(x_n))$ is increasing by hypothesis
and order-bounded by $f(x) + g(x)$. By $\omega$-continuity of $g$
and Lemma~\ref{lem:add-scale-cont},
$g(x) = \sup_n g(x_n) \cle f(x) + \sup_n (g(x_n) - f(x_n))$,
and hence $g(x) - f(x) \cle \sup_n (g(x_n) - f(x_n))$. The
converse inequality holds by monotonicity of $g - f$.
\end{proof}

\begin{lemma}[Operator-norm boundedness]
\label{lem:linmap-norm}
\uses{def:cone}
\rocqok\rocq{Icones.cones.basic_lemmas.linmap_bounded}
\rocq{Icones.cones.basic_lemmas.linmap_norm}
If $f : P \to Q$ is linear (not necessarily $\omega$-continuous)
then $f(B_P)$ is bounded in $Q$, and we may define
\[
  \cnorm{f} \;=\; \sup_{x \in B_P} \cnorm{f(x)} \;\in\; \Rnn.
\]
\end{lemma}

\begin{proof}
If $f(B_P)$ were unbounded, pick $x_n \in B_P$ with
$\cnorm{f(x_n)} \geq 4^n$. The sum $y_n = \sum_{k=1}^{n}
\tfrac{1}{2^k} x_k$ lies in $B_P$ (by (Normt) and (Normh)) and
the chain $(y_n)$ is increasing, hence has a sup $y \in B_P$ by
(Normc). Linearity and (Normp) give
$\cnorm{f(y)} \geq \cnorm{f(\tfrac{1}{2^n} x_n)} = \tfrac{1}{2^n}
\cnorm{f(x_n)} \geq 2^n$ for every $n$, contradicting
$\cnorm{f(y)} \in \Rnn$. Once $f(B_P)$ is bounded, the supremum
$\cnorm{f}$ exists in $\R$ by Dedekind completeness; in Rocq it
is extracted via $\rocq{Icones.cones.basic_lemmas.linmap_norm}$
(based on $\mathsf{xchoose}$ applied to the boundedness witness).
\end{proof}

\begin{remark}
Strikingly, the boundedness of $f(B_P)$ holds without any
continuity assumption on $f$. The Rocq development records this as
$\rocq{Icones.cones.basic_lemmas.linmap_bounded}$, and exposes the
sup as $\rocq{Icones.cones.basic_lemmas.linmap_norm}$ with bound
witness $\rocq{Icones.cones.basic_lemmas.linmap_norm_ub}$.
\end{remark}

%%%%%%%%%%%%%%%%%%%%%%%%%%%%%%%%%%%%%%%%%%%%%%%%%%%%%%%%%%%%%%%%%%%%%%%%
\section{The category $\Cones$}\label{sec:cat-cones}

\begin{definition}[The category $\Cones$ --- Paper Definition~2.17]
\label{def:cones-cat}
\uses{def:cone, lem:linmap-norm}
\rocqok\rocq{Icones.cones.cone_cat.cones_hom}
\rocq{Icones.cones.cone_cat.cones_id}
\rocq{Icones.cones.cone_cat.cones_comp}
The category $\Cones$ has cones as objects. A morphism $f : P \to Q$
in $\Cones(P,Q)$ is a function which is linear,
$\omega$-continuous, and satisfies $\cnorm{f} \leq 1$. Identity
and composition are inherited from $\Set$.
\end{definition}

\begin{remark}
In the Rocq encoding the operator-norm condition $\cnorm{f} \leq 1$
is replaced by its pointwise reformulation
\[
  \forall x \in P,\quad \cnorm{f(x)}_Q \;\leq\; \cnorm{x}_P,
\]
which is equivalent under Lemma~\ref{lem:linmap-norm} but avoids
having to materialise the operator norm at the level of every
morphism. See the design-choices section below for why this trade
is preferable. The morphism record is
$\rocq{Icones.cones.cone_cat.cones_hom}$, with identity and
composition $\rocq{Icones.cones.cone_cat.cones_id}$ and
$\rocq{Icones.cones.cone_cat.cones_comp}$. Category laws are proved
by subtype extensionality ($\rocq{Icones.cones.cone_cat.cones_hom_eq}$)
followed by direct calculation:
$\rocq{Icones.cones.cone_cat.cones_compIl}$,
$\rocq{Icones.cones.cone_cat.cones_compIr}$,
$\rocq{Icones.cones.cone_cat.cones_compA}$.
\end{remark}

\begin{theorem}[Small products in $\Cones$ --- Paper Theorem~2.18]
\label{thm:cones-products}
\uses{def:cones-cat}
\rocqok\rocq{Icones.cones.cone_cat.cones_prod}
\rocq{Icones.cones.cone_cat.cones_proj}
\rocq{Icones.cones.cone_cat.cones_tuple}
Let $(P_i)_{i \in I}$ be a family of cones, indexed by an arbitrary
type $I$. Their product in $\Cones$ is the cone $\&_{i \in I} P_i$
of norm-bounded sections of $\prod_i P_i$, equipped with
componentwise operations and norm
$\cnorm{\vec{x}} = \sup_{i \in I} \cnorm{x_i}_{P_i}$. The
projections are the standard set-theoretic projections; the
mediating arrow associated with $(f_i \in \Cones(Q, P_i))_{i \in I}$
is $y \mapsto (f_i(y))_{i \in I}$.
\end{theorem}

\begin{proof}[Proof sketch]
The carrier is the record
$\rocq{Icones.cones.cone_cat.cones_prod_car}$, a $\Sigma$-type of a
section $i \mapsto x_i$ together with a boundedness witness
$\exists M, \forall i, \cnorm{x_i} \leq M$. Pointwise addition and
scaling preserve boundedness by (Normt) and (Normh); the precone
axioms transport componentwise. The norm
$\cnorm{\vec{x}} = \sup_i \cnorm{x_i}$ is taken in $\R$ via
$\rocq{Icones.cones.cone_cat.cones_prod_norm}$, using Dedekind
completeness on the bounded set
$\rocq{Icones.cones.cone_cat.cones_prod_normset}$. Axioms (Normh)
to (Normp) reduce to componentwise properties; (Normc) is
established by taking componentwise suprema of the chain
($\rocq{Icones.cones.cone_cat.cones_prod_sup_ball}$). Projections
are linear, $\omega$-continuous, and pointwise norm-non-expansive
($\rocq{Icones.cones.cone_cat.cones_proj}$); the tuple
$y \mapsto (f_i(y))_{i \in I}$ is norm-bounded by the sup of the
operator norms of the $f_i$'s, all bounded by $\cnorm{y}$, hence is
itself a morphism in $\Cones$
($\rocq{Icones.cones.cone_cat.cones_tuple}$). Uniqueness of the
mediating arrow follows from extensionality
($\rocq{Icones.cones.cone_cat.cones_tuple_unique}$).
\end{proof}

\begin{remark}
The empty product yields the terminal object $\top$ (paper
\S2.3): the zero-dimensional cone with one element.
\end{remark}

\begin{lemma}[Separate $\omega$-continuity implies joint $\omega$-continuity --- Paper Lemma~2.19]
\label{lem:separate-implies-joint}
\uses{thm:cones-products}
\rocqok\rocq{Icones.cones.cone_cat.diagonal_collapse}
Let $P, Q, R$ be cones, $A \subseteq P$ and $B \subseteq Q$ be
$\omega$-closed subsets, and $f : A \times B \to R$ be separately
$\omega$-continuous (continuous in each argument with the other
fixed). Then $f$ is $\omega$-continuous.
\end{lemma}

\begin{proof}[Proof sketch]
For an $\omega$-chain $((x_n, y_n))$ with sup $(\sup_n x_n, \sup_n y_n)$,
separate continuity gives $f(\sup_n x_n, \sup_n y_n) = \sup_n \sup_k
f(x_n, y_k)$. Monotonicity of $f$ collapses the double supremum to
the diagonal $\sup_n f(x_n, y_n)$. The Rocq lemma
$\rocq{Icones.cones.cone_cat.diagonal_collapse}$ packages the
diagonal-collapse step (the substantive content) on
$\omega$-completed chains.
\end{proof}

\begin{theorem}[Binary equalisers in $\Cones$ --- Paper Theorem~2.20]
\label{thm:cones-equalisers}
\uses{def:cones-cat}
\rocqok\rocq{Icones.cones.cone_cat.cones_eq}
\rocq{Icones.cones.cone_cat.cones_eq_incl}
\rocq{Icones.cones.cone_cat.cones_eq_med}
\rocq{Icones.cones.cone_cat.cones_eq_med_factor}
\rocq{Icones.cones.cone_cat.cones_eq_med_unique}
For any parallel pair $f, g \in \Cones(P, Q)$, the equaliser
$E = \{x \in P \mid f(x) = g(x)\}$ inherits a unique cone structure
from $P$, and the inclusion $e : E \hookrightarrow P$ is a
morphism in $\Cones$. Furthermore $e$ reflects the cone order: if
$e(x) \cle e(y)$ in $P$ then $x \cle y$ in $E$. Together with
Theorem~\ref{thm:cones-products}, this shows that $\Cones$ is
\emph{complete}.
\end{theorem}

\begin{proof}[Proof sketch]
Linearity of $f, g$ makes $E$ closed under the algebraic
operations. The norm of $E$ is the restriction of the norm of $P$,
which trivially inherits (Normh), (Normz), (Normt) and (Normp).
For (Normc), an $\omega$-chain in $B_E$ is in particular a chain
in $B_P$ and has a sup $x \in B_P$; $\omega$-continuity of $f$
and $g$ then yields $f(x) = g(x)$, i.e.\ $x \in E$. The order on
$E$ inherited from $P$ coincides with the cone order computed
inside $E$, because if $x_1 \cle_P x_2$ for $x_1, x_2 \in E$
then the difference $x_2 - x_1$ in $P$ lies in $E$ (by linearity
of $f$ and $g$). Universality of the equaliser is direct: given
$h \in \Cones(T, P)$ with $f \circ h = g \circ h$, the function
$h$ factors through $E$ via the canonical mediator
$\rocq{Icones.cones.cone_cat.cones_eq_med}$; uniqueness follows
from injectivity of the inclusion.
\end{proof}

\begin{lemma}[Norm bound on an isomorphism --- Paper Lemma~2.21 + Proposition~2.22]
\label{lem:iso-preserves-norm}
\uses{def:cones-cat, lem:linmap-norm}
\rocqok\rocq{Icones.cones.cone_cat.cones_iso_preserves_norm}
If $f \in \Cones(P, Q)$ is an isomorphism (i.e.\ there exists
$g \in \Cones(Q, P)$ with $g \circ f = \id_P$ and $f \circ g =
\id_Q$), then $f$ is norm-preserving: $\cnorm{f(x)}_Q = \cnorm{x}_P$
for all $x \in P$.
\end{lemma}

\begin{proof}
One direction $\cnorm{f(x)} \leq \cnorm{x}$ is the morphism
condition on $f$. Conversely, $\cnorm{x} = \cnorm{g(f(x))} \leq
\cnorm{f(x)}$ by the morphism condition on $g$ together with
$g \circ f = \id_P$. The Rocq proof
$\rocq{Icones.cones.cone_cat.cones_iso_preserves_norm}$ records
this in pointwise form, which subsumes Proposition~2.22's
$\cnorm{f} = 1$ (when $P \neq 0$) of the paper.
\end{proof}

\begin{lemma}[Transport of cone structure along a bijection --- Paper Lemma~2.23]
\label{lem:cone-transport}
\uses{def:cone}
\rocqok\rocq{Icones.cones.cone_cat.transport_isPrecone}
Let $P$ be a cone, $S$ a set, and $f : P \to S$ a bijection.
There is a unique cone structure on $S$ making $f$ an isomorphism
in $\Cones$.
\end{lemma}

\begin{proof}[Proof sketch]
Transport the operations along $f$: $s + t := f(f^{-1}(s) +
f^{-1}(t))$, $r \cdot s := f(r \cdot f^{-1}(s))$, $\cnorm{s} :=
\cnorm{f^{-1}(s)}$. The cone axioms transport verbatim, and the
structure is forced by the requirement that $f$ be a $\Cones$-iso.
The Rocq construction lives in
$\rocq{Icones.cones.cone_cat.ConesTransport}$, with the precone
mixin assembled by
$\rocq{Icones.cones.cone_cat.transport_isPrecone}$.
\end{proof}

%%%%%%%%%%%%%%%%%%%%%%%%%%%%%%%%%%%%%%%%%%%%%%%%%%%%%%%%%%%%%%%%%%%%%%%%
\section{Examples: $\top$ and $\bot$}\label{sec:cone-examples}

\begin{definition}[The zero-dimensional cone $\top$]
\label{def:cone-top}
\uses{def:cone}
\rocqok\rocq{Icones.cones.examples_cone.ConeBot.T}
The cone $\top$ has a single element. Addition, scaling, the norm
($\equiv 0$), and the $\omega$-sup are all constant. It is the
terminal object of $\Cones$ (cf.\ Theorem~\ref{thm:cones-products}).
In Rocq, the carrier is the one-constructor inductive
$\rocq{Icones.cones.examples_cone.ConeBot.T}$ wrapped in a
\texttt{Module} to avoid name clashes with the dual example.
\end{definition}

\begin{definition}[The one-dimensional cone $\bot = \Rnn$]
\label{def:cone-bot}
\uses{def:cone}
\rocqok\rocq{Icones.cones.examples_cone.ConeOne.T}
The cone $\bot$ has carrier $\Rnn$, with the obvious operations
$r + s$, $\lambda \cdot r$, and norm $\cnorm{r} = r$. Direct
verification of the cone axioms is bundled in
$\rocq{Icones.cones.examples_cone.ConeOne.T}$.
\end{definition}

\begin{remark}[Wrapper records]
Each of $\top$ and $\bot$ is declared as a fresh inductive (or a
fresh \texttt{Definition} of an abbreviation type) wrapped in a
\texttt{Module} so that the \HB{} instances do not clash. This is
a standard \HB{}/Rocq idiom: two distinct cones must have distinct
carrier types in order for HB to attach two distinct
$\rocq{Icones.cones.precone.isPrecone}$ /
$\rocq{Icones.cones.cone.isCone}$ instances. A raw
\texttt{Notation T := unit} (resp.\ $\texttt{T := \{nonneg R\}}$)
would mean both cones share a carrier, so HB would resolve to one
canonical structure only.
\end{remark}

%%%%%%%%%%%%%%%%%%%%%%%%%%%%%%%%%%%%%%%%%%%%%%%%%%%%%%%%%%%%%%%%%%%%%%%%
\section{Design choices}\label{sec:cones-design}

The Rocq formalisation makes several pragmatic design choices that
diverge from the most literal transcription of paper Definition~2.4.
We document them here because they propagate through M2--M5.

\paragraph{The norm is valued in $\R$, not $\mathsf{\{nonneg\;R\}}$.}
The paper's norm is $\cnorm{\cdot} : P \to \Rnn$. The Rocq mixin
$\rocq{Icones.cones.cone.isCone}$ instead chooses
$\cnorm{\cdot} : P \to R$ with $R$ a $\mathsf{realType}$, and adds
the derived lemma $\cnorm{x} \geq 0$
($\rocq{Icones.cones.cone.cone_norm_ge0}$) to recover the
non-negativity statement. Choosing $\mathsf{\{nonneg\;R\}}$ as codomain
would force a forest of $\%\!:\!\mathsf{num}$ coercions in every
arithmetic step of every downstream proof (e.g.\ in
$\rocq{Icones.cones.cone.cone_norm0}$ alone), with no logical
benefit. Keeping the norm in $R$ and proving positivity separately
matches mathcomp's standard treatment of norms on metric and
normed spaces.

\paragraph{No \texttt{Order.POrder} instance on \texttt{Precone}.}
The cone order
$x_1 \cle x_2 \;\Leftrightarrow\; \exists z,\, x_2 = x_1 + z$
is a \texttt{Prop}-valued partial order, and is in general not
decidable. Attempting to register it as
$\rocq{Icones.cones.precone.Precone}$ +
$\texttt{Order.POrder}$ in HB would either force a $\mathsf{bool}$
projection (forbidden by the existential) or carry the entire
order theory along ghost \texttt{Prop} arguments, with constant
unification grief. We therefore expose the order purely as a
predicate $\rocq{Icones.cones.precone.precone_le}$ with its
reflexivity/transitivity/antisymmetry lemmas (and the rewriting
lemmas of $\rocq{Icones.cones.precone.PreconeOrder}$), reserving
HB hierarchies for the algebraic and normed structures.

\paragraph{(Normc) as a concrete \texttt{sup\_ball} operator.}
The paper's (Normc) is an existence statement: every bounded
$\omega$-chain in $B_P$ has a sup. We could have packaged this in
Rocq using a sigma-type $\{x \mid \cdots\}$ from which a sup could
be extracted by $\mathsf{cid}$ on demand. We chose instead to
expose the sup as a concrete operator
$\rocq{Icones.cones.cone.cone_sup_ball}$ taking the chain and the
chain-hypotheses as arguments, plus three characterising lemmas.
This pays dividends in M2 and M3, where many proofs need direct
access to the sup as a function of the chain (e.g.\ to argue that
two chain-equivalences yield equal sups, by
$\rocq{Icones.prelude.omegacpo.sup_irrelevant}$-style reasoning).
Re-extracting the sup via $\mathsf{cid}$ at every use would block
the rewriting steps in downstream
$\omega$-continuity proofs.

\paragraph{Pointwise vs.\ operator-norm condition on morphisms.}
Paper Definition~2.17 takes morphisms $f : P \to Q$ satisfying
$\cnorm{f} \leq 1$ where $\cnorm{f} = \sup_{x \in B_P} \cnorm{f(x)}$
(Lemma~\ref{lem:linmap-norm}). The Rocq record
$\rocq{Icones.cones.cone_cat.cones_hom}$ instead carries the
pointwise condition
\[
  \forall x \in P, \quad \cnorm{f(x)}_Q \;\leq\; \cnorm{x}_P,
\]
which is logically equivalent (one direction is immediate by
(Normh); the other is by definition of the sup applied to a
rescaling of $x$). The pointwise form has two advantages: it does
not require materialising
$\rocq{Icones.cones.basic_lemmas.linmap_norm}$ inside every
$\mathsf{cones\_hom}$ record (and $\mathsf{linmap\_norm}$ uses
$\mathsf{xchoose}$ internally), and it composes elegantly under
$\rocq{Icones.cones.cone_cat.comp_norm_le1}$. The operator norm
remains available as a derived quantity when needed.
